\documentclass[sigconf]{acmart}

\setcopyright{none}
\copyrightyear{2026}
\acmYear{2026}
\acmDOI{TBD}
\acmConference[Conference Acronym '26]{Conference Name}{Month Day--Day, 2026}{City, Country}
\acmBooktitle{Conference Name (Conference Acronym '26), Month Day--Day, 2026, City, Country}
\acmPrice{15.00}
\acmISBN{978-x-xxxx-xxxx-x/YY/MM}

\usepackage{booktabs}
\usepackage{amsmath}
\usepackage{graphicx}

\title{Practical Patterns for Deterministic AI-Assisted Writing}
\author{Anonymous Research Team}
\affiliation{
  \institution{Institution Name}
  \city{City}
  \country{Country}
}
\email{name@example.com}

\begin{document}

\begin{abstract}
This ACM-style sample presents practical patterns for deterministic AI-assisted manuscript development.
The document shows how structured templates, stable citation keys, and explicit verification checks improve reliability.
It is intended as a shareable baseline for teams adopting the \texttt{codex-latex} skill.
\end{abstract}

\maketitle

\section{Introduction}
AI writing pipelines are effective only when they are auditable and easy to revise.
The sample focuses on simple conventions that reduce formatting drift and reference breakage.

\section{Method}
Our process combines template scaffolding, constrained editing, and diagnostic checks.
Each revision is accepted only after compile and cross-reference validation.

\section{Evaluation}
We summarize outcomes using quality gates rather than model-centric metrics.
Table~\ref{tab:acm-gates} provides an example report format.

\begin{table}[t]
  \centering
  \caption{Illustrative quality-gate outcomes for the workflow.}
  \label{tab:acm-gates}
  \begin{tabular}{lc}
    \toprule
    Gate & Status \\
    \midrule
    Compile pass & Yes \\
    Citation integrity & Yes \\
    Label integrity & Yes \\
    \bottomrule
  \end{tabular}
\end{table}

\section{Discussion and Limitations}
This sample is intentionally concise and does not include dataset-specific details.
Teams should replace placeholders with venue-required metadata and empirically grounded results.

\section{Conclusion}
The generated ACM sample is suitable for sharing as a practical starter kit.
It preserves structure and can be extended without changing the core workflow.

\bibliographystyle{ACM-Reference-Format}
\bibliography{references}

\end{document}
